\section{Banco de questões — Direito Constitucional}

\subsection*{N1 — Fundamento competitivo}
\begin{questao}{\nivel{N1} 05.01 — Cebraspe/Certo ou Errado}A Constituição de 1988 é classificada, tradicionalmente, como escrita, promulgada, rígida e analítica.\end{questao}
\begin{questao}{\nivel{N1} 05.02 — Cebraspe/Certo ou Errado}Soberania é atributo de cada Estado-membro da federação brasileira, enquanto a República Federativa do Brasil possui apenas autonomia.\end{questao}
\begin{questao}{\nivel{N1} 05.03 — Cebraspe/Certo ou Errado}Direitos fundamentais não são, em regra, absolutos e podem entrar em colisão com outros direitos e valores constitucionais.\end{questao}
\begin{questao}{\nivel{N1} 05.04 — Cebraspe/Certo ou Errado}Nacionalidade e cidadania são conceitos equivalentes, pois todo nacional brasileiro exerce necessariamente direitos políticos.\end{questao}
\begin{questao}{\nivel{N1} 05.05 — Cebraspe/Certo ou Errado}A União, os Estados, o Distrito Federal e os Municípios são entes autônomos da federação, observadas as competências constitucionais.\end{questao}
\begin{questao}{\nivel{N1} 05.06 — Cebraspe/Certo ou Errado}Os Tribunais de Contas integram o Poder Judiciário por exercerem julgamento de contas.\end{questao}
\begin{questao}{\nivel{N1} 05.07 — Cebraspe/Certo ou Errado}Lei complementar e lei ordinária se distinguem, entre outros aspectos, pela matéria constitucionalmente reservada e pelo quórum de aprovação.\end{questao}
\begin{questao}{\nivel{N1} 05.08 — Cebraspe/Certo ou Errado}Toda lei complementar é hierarquicamente superior a qualquer lei ordinária, independentemente da matéria tratada.\end{questao}
\begin{questao}{\nivel{N1} 05.09 — Cebraspe/Certo ou Errado}O CNJ integra a estrutura do Poder Judiciário, mas não exerce função jurisdicional típica de tribunal revisor das decisões judiciais.\end{questao}
\begin{questao}{\nivel{N1} 05.10 — Cebraspe/Certo ou Errado}PPA, LDO e LOA integram o sistema constitucional de planejamento e orçamento público.\end{questao}

\subsection*{N2 — Aplicação}
\begin{questao}{\nivel{N2} 05.11 — FGV/princípios fundamentais}Assinale a alternativa que apresenta fundamento da República Federativa do Brasil.\begin{enumerate}[label=\Alph*)]\item Erradicação da pobreza.\item Dignidade da pessoa humana.\item Defesa da paz.\item Cooperação entre os povos.\item Desenvolvimento nacional.\end{enumerate}\end{questao}
\begin{questao}{\nivel{N2} 05.12 — FCC/direitos fundamentais}Uma lei restringe direito fundamental para proteger outro valor constitucional. A análise adequada exige:\begin{enumerate}[label=\Alph*)]\item considerar todo direito fundamental absoluto;\item examinar fundamento constitucional, proporcionalidade e preservação do núcleo protegido;\item presumir a lei inválida por existir restrição;\item ignorar a finalidade da medida;\item avaliar apenas conveniência política.\end{enumerate}\end{questao}
\begin{questao}{\nivel{N2} 05.13 — Cebraspe/Certo ou Errado}O habeas data pode ser instrumento destinado ao conhecimento ou à retificação de informações relativas ao próprio impetrante em registros nos termos constitucionais e legais.\end{questao}
\begin{questao}{\nivel{N2} 05.14 — FGV/remédios constitucionais}A ausência de norma regulamentadora inviabiliza o exercício de prerrogativa constitucional. O remédio constitucional especificamente vocacionado ao problema é:\begin{enumerate}[label=\Alph*)]\item habeas corpus;\item ação popular;\item mandado de injunção;\item habeas data;\item mandado de segurança coletivo necessariamente.\end{enumerate}\end{questao}
\begin{questao}{\nivel{N2} 05.15 — FCC/federação}A competência legislativa concorrente caracteriza-se, em linhas gerais, por:\begin{enumerate}[label=\Alph*)]\item exclusividade absoluta dos Municípios;\item repartição constitucional entre União, Estados e DF nos termos do art. 24;\item inexistência de normas gerais federais;\item competência residual da União;\item competência administrativa comum apenas.\end{enumerate}\end{questao}
\begin{questao}{\nivel{N2} 05.16 — Cebraspe/Certo ou Errado}A intervenção federal é medida excepcional e somente pode ocorrer nas hipóteses previstas na Constituição.\end{questao}
\begin{questao}{\nivel{N2} 05.17 — FGV/Administração Pública}A investidura em cargo efetivo depende, em regra, de:\begin{enumerate}[label=\Alph*)]\item livre nomeação;\item concurso público prévio, observadas as exceções constitucionais;\item eleição direta;\item indicação do Tribunal de Contas;\item contrato administrativo.\end{enumerate}\end{questao}
\begin{questao}{\nivel{N2} 05.18 — FCC/acumulação}A acumulação remunerada de cargos públicos é:\begin{enumerate}[label=\Alph*)]\item sempre permitida;\item vedada como regra, com exceções constitucionais condicionadas aos requisitos previstos;\item permitida apenas para agentes políticos;\item proibida inclusive nas hipóteses constitucionais;\item matéria exclusivamente infraconstitucional.\end{enumerate}\end{questao}
\begin{questao}{\nivel{N2} 05.19 — Cebraspe/Certo ou Errado}O servidor estável pode perder o cargo em hipóteses constitucionalmente previstas; estabilidade não significa impossibilidade absoluta de desligamento.\end{questao}
\begin{questao}{\nivel{N2} 05.20 — FGV/processo legislativo}Uma proposta de emenda constitucional pretende abolir a separação dos Poderes. A proposta:\begin{enumerate}[label=\Alph*)]\item pode ser aprovada por maioria simples;\item encontra limitação material decorrente de cláusula pétrea;\item exige apenas sanção presidencial;\item pode ser aprovada por medida provisória;\item depende apenas de referendo.\end{enumerate}\end{questao}
\begin{questao}{\nivel{N2} 05.21 — FCC/controle externo}No plano federal, a fiscalização contábil, financeira e orçamentária é exercida pelo Congresso Nacional mediante controle externo com auxílio:\begin{enumerate}[label=\Alph*)]\item do STF;\item do TCU;\item do CNJ;\item do Ministério Público Federal exclusivamente;\item da CGU apenas.\end{enumerate}\end{questao}
\begin{questao}{\nivel{N2} 05.22 — Cebraspe/Certo ou Errado}O TCU aprecia as contas prestadas anualmente pelo Presidente da República, emitindo parecer prévio, enquanto julga contas de administradores e demais responsáveis sujeitos à sua jurisdição.\end{questao}
\begin{questao}{\nivel{N2} 05.23 — FGV/Executivo}A medida provisória:\begin{enumerate}[label=\Alph*)]\item depende de relevância e urgência e está sujeita a limitações constitucionais;\item é lei complementar;\item pode tratar qualquer matéria sem exceção;\item independe de apreciação do Congresso;\item somente existe nos Municípios.\end{enumerate}\end{questao}
\begin{questao}{\nivel{N2} 05.24 — FCC/Judiciário}O Conselho Nacional de Justiça exerce, principalmente:\begin{enumerate}[label=\Alph*)]\item revisão jurisdicional de sentenças;\item controle administrativo e financeiro do Judiciário e deveres funcionais, nos termos constitucionais;\item controle externo do Executivo;\item julgamento de contas públicas;\item processo legislativo.\end{enumerate}\end{questao}
\begin{questao}{\nivel{N2} 05.25 — Cebraspe/Certo ou Errado}A ação popular pode ser proposta por cidadão nas hipóteses constitucionais e legais para proteger bens e valores indicados na Constituição.\end{questao}

\subsection*{N3 — Integração de concurso}
\begin{questao}{\nivel{N3} 05.26 — FGV/assertivas direitos}Considere: I. Direitos fundamentais têm aplicação imediata. II. Aplicação imediata significa que nenhum direito depende jamais de regulamentação. III. Direitos fundamentais podem sofrer restrições constitucionalmente justificadas. Assinale:\begin{enumerate}[label=\Alph*)]\item apenas I;\item apenas II;\item I e III;\item II e III;\item I, II e III.\end{enumerate}\end{questao}
\begin{questao}{\nivel{N3} 05.27 — Cebraspe/Certo ou Errado}A igualdade constitucional não impede tratamentos diferenciados quando exista fundamento constitucionalmente legítimo e critério proporcional à diferença relevante.\end{questao}
\begin{questao}{\nivel{N3} 05.28 — FCC/federação}Matéria de interesse predominantemente local é atribuída constitucionalmente, em regra, aos:\begin{enumerate}[label=\Alph*)]\item Municípios;\item Estados exclusivamente;\item Senado;\item Tribunais de Contas;\item partidos políticos.\end{enumerate}\end{questao}
\begin{questao}{\nivel{N3} 05.29 — FGV/competências}Os Estados possuem, em regra, competências:\begin{enumerate}[label=\Alph*)]\item apenas expressamente enumeradas;\item remanescentes/residuais, ressalvadas vedações constitucionais;\item exclusivas da União;\item idênticas às dos Municípios;\item apenas delegadas pelo Presidente.\end{enumerate}\end{questao}
\begin{questao}{\nivel{N3} 05.30 — Cebraspe/Certo ou Errado}A competência comum do art. 23 é predominantemente material/administrativa e não deve ser confundida com a competência legislativa concorrente do art. 24.\end{questao}
\begin{questao}{\nivel{N3} 05.31 — FCC/servidores}A estabilidade do servidor efetivo:\begin{enumerate}[label=\Alph*)]\item equivale à vitaliciedade;\item é adquirida segundo os requisitos constitucionais e não se estende automaticamente a ocupante exclusivamente comissionado;\item impede avaliação de desempenho;\item nasce antes do concurso;\item independe de cargo efetivo.\end{enumerate}\end{questao}
\begin{questao}{\nivel{N3} 05.32 — FGV/processo legislativo}Lei ordinária é aprovada por maioria simples, observada a presença necessária; lei complementar exige:\begin{enumerate}[label=\Alph*)]\item unanimidade;\item maioria absoluta;\item três quintos em dois turnos;\item dois terços;\item referendo.\end{enumerate}\end{questao}
\begin{questao}{\nivel{N3} 05.33 — Cebraspe/Certo ou Errado}A iniciativa reservada de determinada matéria integra o controle formal de constitucionalidade da lei.\end{questao}
\begin{questao}{\nivel{N3} 05.34 — FCC/Tribunais de Contas}Assinale a alternativa correta.\begin{enumerate}[label=\Alph*)]\item O Tribunal de Contas é subordinado hierarquicamente ao Legislativo.\item O Tribunal de Contas auxilia o Legislativo e possui competências constitucionais próprias.\item O Tribunal de Contas integra o Judiciário.\item O Tribunal de Contas apenas emite pareceres e nunca aplica sanções.\item O Tribunal de Contas substitui o controle interno.\end{enumerate}\end{questao}
\begin{questao}{\nivel{N3} 05.35 — FGV/contas}Em relação às contas de governo do chefe do Executivo e às contas de gestão dos administradores, a Constituição federal estabelece, como paradigma:\begin{enumerate}[label=\Alph*)]\item julgamento de ambas exclusivamente pelo TCU;\item parecer prévio sobre as contas do Presidente e julgamento das contas dos administradores/responsáveis pelo TCU;\item parecer prévio em todas as contas sem julgamento;\item julgamento de todas pelo STF;\item inexistência de distinção.\end{enumerate}\end{questao}
\begin{questao}{\nivel{N3} 05.36 — Cebraspe/Certo ou Errado}A fiscalização constitucional abrange legalidade, legitimidade e economicidade, razão pela qual controle externo não se limita à conformidade formal da despesa.\end{questao}
\begin{questao}{\nivel{N3} 05.37 — FCC/controle de constitucionalidade}No controle difuso, a questão constitucional pode ser examinada:\begin{enumerate}[label=\Alph*)]\item apenas pelo STF em ação direta;\item incidentalmente em caso concreto por órgãos jurisdicionais competentes;\item somente pelo Senado;\item pelo Tribunal de Contas como ação direta;\item exclusivamente pelo Presidente.\end{enumerate}\end{questao}
\begin{questao}{\nivel{N3} 05.38 — FGV/controle concentrado}A ADI de lei federal ou estadual em face da Constituição Federal, nas hipóteses constitucionais, insere-se no controle:\begin{enumerate}[label=\Alph*)]\item difuso municipal;\item concentrado perante o STF;\item interno administrativo;\item legislativo preventivo apenas;\item financeiro.\end{enumerate}\end{questao}
\begin{questao}{\nivel{N3} 05.39 — Cebraspe/Certo ou Errado}Os efeitos das decisões em controle de constitucionalidade não podem ser reduzidos a uma regra absoluta ``difuso sempre inter partes e concentrado sempre sem modulação'', pois o sistema possui técnicas e efeitos mais complexos.\end{questao}
\begin{questao}{\nivel{N3} 05.40 — FCC/orçamento}A lei que estima receitas e fixa despesas para o exercício é a:\begin{enumerate}[label=\Alph*)]\item LDO;\item LOA;\item PPA;\item lei delegada;\item resolução do TCU.\end{enumerate}\end{questao}

\subsection*{N4 — Auditoria / avançado}
\begin{questao}{\nivel{N4} 05.41 — Cebraspe/caso integrado}O Legislativo determina ao Tribunal de Contas o resultado que deverá adotar em julgamento de contas específico, invocando relação hierárquica. Julgue: a premissa é incompatível com a autonomia funcional e as competências constitucionais próprias do Tribunal de Contas.\end{questao}
\begin{questao}{\nivel{N4} 05.42 — FGV/controle externo}Auditoria identifica despesa formalmente prevista, mas com preço muito superior ao mercado e sem justificativa. A avaliação de economicidade pelo Tribunal de Contas:\begin{enumerate}[label=\Alph*)]\item é vedada porque só legalidade pode ser examinada;\item integra o alcance constitucional da fiscalização;\item compete apenas ao Judiciário;\item exige crime previamente reconhecido;\item não se aplica a contratos.\end{enumerate}\end{questao}
\begin{questao}{\nivel{N4} 05.43 — FCC/direitos fundamentais}Uma política pública utiliza tratamento distinto entre grupos. Para avaliar constitucionalidade, é necessário:\begin{enumerate}[label=\Alph*)]\item considerar qualquer distinção proibida;\item verificar finalidade legítima, critério de diferenciação e proporcionalidade;\item ignorar igualdade;\item considerar apenas custo;\item presumir discriminação ilícita em toda classificação.\end{enumerate}\end{questao}
\begin{questao}{\nivel{N4} 05.44 — Cebraspe/Certo ou Errado}O devido processo constitucional pode exigir oportunidade efetiva de manifestação antes de decisão estatal gravosa, não bastando comunicação meramente formal quando ela inviabiliza defesa.\end{questao}
\begin{questao}{\nivel{N4} 05.45 — FGV/federação}A União pretende legislar sobre matéria de interesse local sem fundamento em competência federal. O primeiro ponto de análise constitucional é:\begin{enumerate}[label=\Alph*)]\item a repartição de competências federativas;\item a responsabilidade civil do Estado;\item a vitaliciedade;\item a neutralidade da Internet;\item a lei orçamentária.\end{enumerate}\end{questao}
\begin{questao}{\nivel{N4} 05.46 — Cebraspe/Certo ou Errado}A autonomia dos entes federativos envolve capacidades de auto-organização, autogoverno, autoadministração e autolegislação nos limites constitucionais, mas não soberania própria.\end{questao}
\begin{questao}{\nivel{N4} 05.47 — FCC/processo legislativo}Projeto de lei ordinária trata matéria expressamente reservada pela Constituição à lei complementar. O vício é principalmente:\begin{enumerate}[label=\Alph*)]\item material de finalidade;\item formal, por inadequação da espécie normativa/procedimento;\item inexistente;\item administrativo apenas;\item orçamentário.\end{enumerate}\end{questao}
\begin{questao}{\nivel{N4} 05.48 — FGV/TC e sanções}Ao constatar ilegalidade em ato sujeito à sua fiscalização, o Tribunal de Contas pode exercer competências constitucionais de determinação/correção e sanção nos limites aplicáveis. Isso demonstra que:\begin{enumerate}[label=\Alph*)]\item sua atuação é exclusivamente opinativa;\item possui competências decisórias próprias além do parecer prévio sobre certas contas;\item integra o Executivo;\item substitui o Ministério Público;\item não exerce fiscalização operacional.\end{enumerate}\end{questao}
\begin{questao}{\nivel{N4} 05.49 — Cebraspe/Certo ou Errado}A existência de cláusulas pétreas limita o poder constituinte derivado reformador, embora não impeça toda e qualquer alteração textual relacionada a temas protegidos.\end{questao}
\begin{questao}{\nivel{N4} 05.50 — FGV/matriz de análise}Uma lei estadual invade competência legislativa privativa da União. A análise de constitucionalidade deve verificar prioritariamente:\begin{enumerate}[label=\Alph*)]\item competência legislativa e eventual autorização constitucional de delegação;\item eficiência administrativa;\item mérito político;\item conveniência do governador;\item apenas impacto financeiro.\end{enumerate}\end{questao}
